\section{Cái chương này đóng được, và cái nó mở ra}
\label{sec:ch09-con-lai-gi}

\index{tiền huấn luyện!chương này đóng được gì}%
Chương~\ref{ch:chi-phi-va-song-song} để lại ba khoảng trống của kiến trúc:
chi phí bậc hai theo độ dài, không có \tn{thiên kiến quy nạp}{inductive bias}
về vị trí, và đói dữ liệu. Chương ấy đo được khoảng trống thứ nhất và loại trừ được khoảng trống thứ
hai khỏi danh sách nghi phạm. Chương này là chương của khoảng trống thứ ba.

Cái nó đóng được thì gọn, và không phải cái tôi định đóng khi bắt đầu.

Kiến trúc không đổi. Ba bài mở mục kiến trúc bằng ba câu nói rằng chúng dùng
lại khối của Vaswani, và hình~\ref{fig:ch09-ba-cach-cat} nói lại chuyện ấy bằng
hình: encoder-only và decoder-only là hai phép xóa trên cùng một khối, và cái
phân biệt hai phép xóa còn lại với nhau là một ma trận tam giác.

Cái đổi là mục tiêu huấn luyện, và mục tiêu ấy chọn được vì nó không cần nhãn.
Mọi thứ còn lại theo sau: bỏ mask thì mất mục tiêu đoán token tiếp theo, nên
phải có masked language model; masked language model chấm 15\% số vị trí nên
đắt hơn 6.67 lần cho mỗi đích; và cả hai đều đáng, vì pha một trả một lần còn
pha hai trả cho từng bài toán.

\subsection*{Ba con số không quay về, và một cái bảng đếm thiếu}

\index{tiền huấn luyện!bốn chỗ dựng lại không khớp}%
Số tham số của BERT-large in 340M và dựng lại ra 335.14M, lệch 1.4\%, không
đóng được dưới bất kỳ cách đọc nào tôi thử. Bảng 2 của GPT-2 in bốn con số chỉ
tái lập được bằng một từ điển quanh 40,500, trong khi mục 2.3 cùng bài thông báo
từ điển đã mở rộng lên 50,257. GPT-3 in 175.0B ở bảng 2.1 và 174,600M ở bảng
D.1 của chính nó, và phép dựng lại cho 174,604,259,328.

GPT-1 thì không có con số nào để đối chiếu, vì bài không in số tham số.

Không chỗ nào trong ba chỗ ấy làm đổi một luận điểm nào. Chúng đáng in ra vì lý
do khác, và lý do ấy là cái tôi mang ra khỏi phiên làm việc này: \textbf{một con
số được trích lại đủ nhiều lần thì không còn ai dựng lại nó.} Phép dựng lại ở
đây là số học dạng đóng trên siêu tham số các bài tự in, không huấn luyện gì và
không đo đồng hồ.

Bảng 1 của Kaplan cũng vậy. Nó đếm một trong hai tích bậc hai của attention,
tỷ lệ đúng 2.00 ở mọi cấu hình, và chuyện ấy không đi tới đâu vì chính số hạng
đó bị \(6N\) vứt đi ở dòng sau. Nhưng ai dựng phép đếm FLOP của mình theo bảng
ấy sẽ mang theo chỗ thiếu, và bảng ấy là chỗ bài đặt công thức sẵn để chép.

\subsection*{Chỗ hai chương nối vào nhau}

\index{tiền huấn luyện!ngưỡng 6d nối hai chương}%
Một chuyện tôi không chờ đợi: sai số của \(6ND\) và phát hiện của
chương~\ref{ch:chi-phi-va-song-song} là cùng một đại lượng.

\eqref{eq:ch09-phan-attention} nói phần attention của một lượt xuôi là
\(n/(6d+n)\), bằng nửa tại \(n = 6d\). Mục~\ref{sec:ch08-dem} tìm ra
\(n = 2d + \dff\) là chỗ phần bậc hai vượt nửa một tầng, và ở \(\dff = 4d\) thì
đó đúng là \(6d\). Chương trước hỏi khi nào phần bậc hai đắt bằng phần còn lại;
chương này hỏi khi nào xấp xỉ mà cả ngành dùng để lập ngân sách bỏ sót một nửa.
Cùng một ngưỡng.

Hệ quả đọc được ngay từ bảng của mục~\ref{sec:ch09-sau-nd}: \(6ND\) sai 2.8\%
ở GPT-3 và 44\% ở đúng mô hình ấy với ngữ cảnh 32768. Không phải vì quy mô. Vì
\(n/d\). Mọi thứ trong hai bài \tn{quy luật scaling}{scaling law} đứng trên một
xấp xỉ đúng trong chế độ \(n \ll d\), và chế độ ấy là chế độ của 2020.

\subsection*{Cái nó mở ra}

\index{tiền huấn luyện!chỗ ba chương sau bắt đầu}%
Ba chỗ, và chúng đi về ba hướng khác nhau.

Thứ nhất, chương~\ref{ch:bias-quy-nap} và chương~\ref{ch:vit}. Chương này lập
xong tiền đề mà Vision Transformer, viết tắt là ViT, cần: cùng một kiến trúc
thắng hay thua tùy vào nó được \tn{tiền huấn luyện}{pretraining} trên cái gì.
Không có tiền đề ấy thì kết quả trung tâm của ViT đọc
như một nghịch lý, vì nó là một bảng trong đó cùng một mô hình đứng hai bên của
một phép so.

Thứ hai, chương~\ref{ch:vi-tri-va-ngu-canh}. Điều kiện của Kaplan,
\(d_{\mathrm{model}} \gg n_{\mathrm{ctx}}/12\), là điều kiện của các mô hình
2020 và không phải điều kiện của mô hình ngữ cảnh dài. Bảng ở
mục~\ref{sec:ch09-sau-nd} đo giá của việc bỏ điều kiện ấy: cùng GPT-3, ngữ cảnh
gấp mười sáu, sai số từ 2.8\% lên 44\%. Mọi ngân sách tính bằng \(6ND\) cho một
mô hình ngữ cảnh dài đều thấp, và thấp một khoảng tính được.

Thứ ba, và đây là chỗ mở thật sự. Hai bài trong chương này bất đồng về một câu
hỏi mà cả hai đều trả lời bằng phép khớp thực nghiệm, và bài sau chỉ ra rằng
bài trước sai vì một chi tiết của lịch learning rate. Không ai biết bài sau còn
bao nhiêu chi tiết như thế trong nó. Bảng A4 của chính bài ấy đã không dựng lại
được từ những gì bài in, và bảng 2 của nó không dựng lại được từ hằng số phụ lục
D.2 của nó. Bài tập tier ba cuối chương giao lại phần ấy, và tôi ghi ở đây rằng
nó chưa được giải chứ không ghi rằng nó là bài tập.
